% BHU PYQ -- Professional Exam Template
\documentclass[11pt,a4paper]{article}

\usepackage[a4paper,top=2.5cm,bottom=2.5cm,left=2.8cm,right=2.8cm,headheight=14pt]{geometry}
\usepackage{amsmath,amssymb}
\usepackage[T1]{fontenc}
\usepackage[utf8]{inputenc}
\usepackage{lmodern}
\usepackage{microtype}
\usepackage{fancyhdr}
\usepackage{enumitem}
\usepackage{listings}
\usepackage{hyperref}
\usepackage{lastpage}
\usepackage{parskip}

\hypersetup{colorlinks=true,urlcolor=black,linkcolor=black,
  pdftitle={BVCA-304: Java Applet Programming -- BHU PYQ}}

\pagestyle{fancy}
\fancyhf{}
\renewcommand{\headrulewidth}{0.4pt}
\renewcommand{\footrulewidth}{0.4pt}
\fancyhead[L]{\small\textit{Banaras Hindu University}}
\fancyhead[R]{\small\textit{BVCA-304: Java Applet Programming}}
\fancyfoot[C]{\small Page \thepage\ of \pageref{LastPage}}

\lstset{
  basicstyle=\small\ttfamily,
  language=Java,
  frame=single,
  breaklines=true,
  columns=flexible,
  keepspaces=true
}

\newlist{parts}{enumerate}{2}
\setlist[parts,1]{label=(\alph*),leftmargin=*,itemsep=4pt,topsep=3pt}
\setlist[parts,2]{label=(\roman*),leftmargin=*,itemsep=2pt,topsep=2pt}
\newcommand{\pts}[1]{\hfill\textit{[#1]}}

\begin{document}

\begin{center}
  {\large\bfseries BANARAS HINDU UNIVERSITY}\\[2pt]
  {\normalsize B.Voc. Semester III Examination, 2022-23}\\[6pt]
  \rule{0.8\linewidth}{0.5pt}\\[6pt]
  {\large\bfseries Paper: BVCA-304 --- Java Applet Programming}\\[4pt]
  {\normalsize Time: Three Hours \hfill Maximum Marks: 70}
\end{center}

\rule{\linewidth}{0.4pt}

\smallskip
\noindent\textit{Instructions: Answer all sections. Write your answers in your own words as far as practicable.}

\bigskip

\bigskip\noindent{\large\bfseries SECTION --- A}
\rule{\linewidth}{0.4pt}\smallskip

\noindent\textit{Note: Long-answer questions to be answered in about 750 words each. Attempt any two.} \hfill \textbf{[15 $\times$ 2 = 30]}

\medskip\noindent\textbf{Question 1.} What are data types and variables in Java? Differentiate between static, instance, and local variables. Why can keywords not be used as variable names?
\centerline{\textbf{OR}}
Discuss the \texttt{if-else} statement and \texttt{switch} statement. Write a program to display the day of the week for a given number using a \texttt{switch} statement. (Days: Sunday=0, Monday=1, Tuesday=2, \ldots)

\medskip\noindent\textbf{Question 2.} What are arrays? Explain their types and discuss the ways to initialize them.
\begin{parts}
  \item Write a program to find the maximum element in an array of 20 integer values.
\end{parts}
\centerline{\textbf{OR}}
Explain the \texttt{for}, \texttt{while}, and \texttt{do-while} loop structures in Java with examples. What is the difference between entry-controlled loops and exit-controlled loops?

\bigskip\noindent{\large\bfseries SECTION --- B}
\rule{\linewidth}{0.4pt}\smallskip

\noindent\textit{Note: Short-answer questions to be answered in about 200 words each. Attempt any six.} \hfill \textbf{[5 $\times$ 6 = 30]}

\medskip\noindent\textbf{Question 3.} How do you read data using the \texttt{Scanner} class in Java? State at least three methods defined in the \texttt{Scanner} class to read data.

\medskip\noindent\textbf{Question 4.} What is concatenation of Strings? Write a program to concatenate two strings.

\medskip\noindent\textbf{Question 5.} Differentiate between the \texttt{break} and \texttt{continue} statements with examples.

\medskip\noindent\textbf{Question 6.} What are the features of the Java programming language?

\medskip\noindent\textbf{Question 7.} Write a program to check whether a number is prime or not.

\medskip\noindent\textbf{Question 8.} Define JDK and JVM. Why is Java considered a platform-independent language?

\medskip\noindent\textbf{Question 9.} What is a Constructor? Explain its types.

\medskip\noindent\textbf{Question 10.} What is an Object and how is it created? Create a class \texttt{Student} with at least three attributes and create two objects of that class.

\bigskip\noindent{\large\bfseries SECTION --- C}
\rule{\linewidth}{0.4pt}\smallskip

\noindent\textit{Note: Answer all objective-type questions.} \hfill \textbf{[10 $\times$ 1 = 10]}

\medskip\noindent\textbf{Question 11.} Answer the following:
\begin{enumerate}[label=(\roman*),leftmargin=*,itemsep=8pt]

  \item Who invented the Java programming language?
  \begin{enumerate}[label=(\alph*),itemsep=2pt]
    \item Guido van Rossum
    \item James Gosling
    \item Dennis Ritchie
    \item Bjarne Stroustrup
  \end{enumerate}

  \item \underline{\hspace{2cm}} is used to find and fix bugs in Java programs.
  \begin{enumerate}[label=(\alph*),itemsep=2pt]
    \item JVM
    \item JRE
    \item JDK
    \item JDB
  \end{enumerate}

  \item What does an Eclipse IDE allow?
  \begin{enumerate}[label=(\alph*),itemsep=2pt]
    \item Editing the program only
    \item Compiling the program only
    \item Both Editing and Compiling the program
    \item Editing, Building, Debugging, and Compiling the program
  \end{enumerate}

  \item SDK stands for \underline{\hspace{2cm}}.
  \begin{enumerate}[label=(\alph*),itemsep=2pt]
    \item Software Design Kit
    \item Serial Development Kit
    \item Software Development Kit
    \item Serial Design Kit
  \end{enumerate}

  \item Which of the following is not an OOP concept in Java?
  \begin{enumerate}[label=(\alph*),itemsep=2pt]
    \item Polymorphism
    \item Inheritance
    \item Compilation
    \item Encapsulation
  \end{enumerate}

  \item What will be the output of the following Java program?
\begin{lstlisting}
class Variable_scope {
    public static void main(String args[]) {
        int x;
        x = 5;
        {
            int y = 6;
            System.out.print(x + " " + y);
        }
        System.out.println(x + " " + y);
    }
}
\end{lstlisting}
  \begin{enumerate}[label=(\alph*),itemsep=2pt]
    \item Compilation error
    \item Runtime error
    \item 5 6 5 6
    \item 5 6 5
  \end{enumerate}

  \item What will be the output of the following Java program?
\begin{lstlisting}
class output {
    public static void main(String args[]) {
        StringBuffer s1 = new StringBuffer("Quiz");
        StringBuffer s2 = s1.reverse();
        System.out.println(s2);
    }
}
\end{lstlisting}
  \begin{enumerate}[label=(\alph*),itemsep=2pt]
    \item QuizziuQ
    \item ziuQQuiz
    \item Quiz
    \item ziuQ
  \end{enumerate}

  \item In which memory is a String stored when we create a String using the \texttt{new} operator?
  \begin{enumerate}[label=(\alph*),itemsep=2pt]
    \item Stack
    \item String memory
    \item Heap memory
    \item Random storage space
  \end{enumerate}

  \item Which keyword is used for accessing the features of a package?
  \begin{enumerate}[label=(\alph*),itemsep=2pt]
    \item \texttt{package}
    \item \texttt{import}
    \item \texttt{extends}
    \item \texttt{export}
  \end{enumerate}

  \item Which value is returned by the \texttt{read()} method when End of File (EOF) is encountered?
  \begin{enumerate}[label=(\alph*),itemsep=2pt]
    \item 0
    \item 1
    \item $-1$
    \item Null
  \end{enumerate}

\end{enumerate}

\vfill
\rule{\linewidth}{0.4pt}
\begin{center}\small
\href{https://bhu-exam.vercel.app/}{Click here to visit our website}\\[4pt]\href{https://me-aryan.vercel.app/}{Click here to visit our contributor}\\[4pt]\href{https://quashan-me.vercel.app/}{Click here to visit our contributor}
\end{center}

\end{document}
